\section{Odd-subset lift characterization}
\label{sec:odd-subset-lifts}

The even-cycle condition also has an exact formulation without auxiliary
cycle colours.  This formulation will be useful both structurally and for
proof-producing computation.

\begin{lemma}[Invariant odd subsets]\label{lem:invariant-odd-subset}
A permutation of a finite set has only even cycles if and only if it preserves
no subset of odd cardinality.
\end{lemma}

\begin{proof}
The support of an odd cycle is an invariant subset of odd cardinality.
Conversely, every invariant subset is a disjoint union of cycles.  A union of
even cycles has even cardinality, so an invariant subset of odd cardinality
contains an odd cycle.
\end{proof}

Let $\Lambda=\{\lambda_i:U\to V:i\in I\}$ be one line family of a Latin
square.  Thus, in the row view, $U$ is the column set, $V$ is the symbol set,
and $\lambda_i$ is row $i$ as a bijection.  The column view uses bijections
from rows to symbols, and the symbol view uses the position bijections from
columns to rows.  In every case $|I|=|U|=|V|=n$, and distinct lines are
pointwise disjoint.

For $1\leq k\leq n$, let $\rho_k(\lambda_i)$ be the permutation matrix of the
induced bijection from the $k$-subsets of $U$ to the $k$-subsets of $V$, and
put
\begin{equation}
 B_k(\Lambda)=\sum_{i\in I}\rho_k(\lambda_i).
 \label{eq:odd-subset-lift}
\end{equation}
Its $(A,B)$ entry counts the lines satisfying $\lambda_i(A)=B$.

\begin{theorem}[Odd-subset collision characterization]
\label{thm:odd-subset-collision}
Let $\Lambda$ be a view of a Latin square of even order $n$.  The view is F if
and only if $B_k(\Lambda)$ is a zero-one matrix for every odd
$k$ with $3\leq k\leq n/2$.  If $n/2$ is odd, it is enough at level $n/2$ to
test one member of each complementary pair of subsets.

Consequently, a Latin square of even order is FFF if and only if this
collision-free condition holds for the row, column, and symbol line families.
\end{theorem}

\begin{proof}
For distinct $i,j$ and a $k$-subset $A\subseteq U$, the two induced images
collide precisely when
\[
 \lambda_i(A)=\lambda_j(A)
 \quad\Longleftrightarrow\quad
 (\lambda_j^{-1}\lambda_i)(A)=A.
\]
Thus an entry of $B_k(\Lambda)$ exceeds one precisely when a relative
two-line permutation preserves the corresponding $k$-subset.

By Lemma~\ref{lem:invariant-odd-subset}, a relative permutation has an odd
cycle if and only if it preserves an odd-cardinality subset.  Since two
distinct Latin-square lines have no common point, their relative permutation
is fixed-point-free, so $k=1$ is automatic.  Because $n$ is even, a subset is
invariant if and only if its complement is invariant, and the two sets have
the same odd parity.  It therefore suffices to retain sizes at most $n/2$;
when $k=n/2$ is odd, one set from each complementary pair is enough.  This
proves the claim for one view.  Applying the same argument to the three line
families proves the FFF statement.
\end{proof}

Every matrix $B_k(\Lambda)$ has constant row and column sum $n$.  Hence, in
an F-view, it is the biadjacency matrix of a simple $n$-regular bipartite
graph on the two $k$-subset sets.  At order $10$, the theorem requires exactly
the levels $k=3$ and $k=5$, with one representative of each complementary
pair at level five.  At orders $6$ and $8$, only $k=3$ is needed.

The characterization has a direct CNF consequence.  If $P_{ij}(x,y)$ denotes
the exact permutation matrix of $\lambda_j^{-1}\lambda_i$, then for every
tested odd subset $A$ one may add the clause
\[
 \bigvee_{x\in A,\ y\notin A} P_{ij}(x,y).
\]
Because $P_{ij}$ is a permutation matrix, this clause says exactly that
$P_{ij}(A)\neq A$.  Theorem~\ref{thm:odd-subset-collision} proves that the
resulting three-view encoding is equisatisfiable with FFF.  Formula size alone
does not imply computational superiority; the bounded order-$10$ pilots in
Section~\ref{sec:reproducibility} remain undecided.
